% !TEX root = ../main.tex
\chapter{Exploratory Data Analysis (EDA)}
\label{ch:eda}

\begin{intuition}
Exploratory Data Analysis (EDA) is an analysis philosophy initiated by John Tukey in the 1970s. Rather than testing predefined hypotheses, EDA invites us to ``let the data speak'' by systematically exploring them through numerical summaries and visualizations.
\end{intuition}

% ============================================================
\section{What is EDA?}
\index{EDA}\index{exploratory analysis}

\begin{definition}[Exploratory Data Analysis]
EDA\index{EDA} is an analytical approach that uses visual and quantitative techniques to understand the structure, anomalies, and relationships in a dataset \emph{before} any formal modeling.
\end{definition}

The main objectives of EDA are:
\begin{itemize}
    \item Discover the \textbf{structure} of the data (dimensions, types, missing values).
    \item Identify \textbf{distributions} and \textbf{central tendencies}.
    \item Detect \textbf{anomalies} and \textbf{outliers}.
    \item Explore \textbf{relationships} between variables.
    \item Formulate \textbf{hypotheses} to be tested later.
\end{itemize}

% ============================================================
\section{Systematic EDA Process}

\begin{center}
\begin{tikzpicture}[
    node distance=1.2cm,
    block/.style={rectangle, draw=steelblue!80, fill=steelblue!10, rounded corners, text width=3.8cm, align=center, minimum height=1cm, font=\small},
    arrow/.style={->, thick, >=stealth, color=gray!70}
]
\node[align=center,block] (load) {1. Load data\\{\scriptsize \py{pd.read\_csv}, \py{sns.load\_dataset}}};
\node[align=center,block, below=of load] (shape) {2. Dimensions and types\\{\scriptsize \py{shape}, \py{dtypes}, \py{info()}}};
\node[align=center,block, below=of shape] (missing) {3. Missing values\\{\scriptsize \py{isnull().sum()}}};
\node[align=center,block, below=of missing] (desc) {4. Descriptive statistics\\{\scriptsize \py{describe()}, \py{value\_counts()}}};
\node[align=center,block, right=2.5cm of load] (uni) {5. Univariate analysis\\{\scriptsize histograms, boxplots}};
\node[align=center,block, below=of uni] (bi) {6. Bivariate analysis\\{\scriptsize scatter, correlation}};
\node[align=center,block, below=of bi] (multi) {7. Multivariate analysis\\{\scriptsize pairplot, heatmap}};
\node[block, below=of multi] (concl) {8. Conclusions and hypotheses};

\draw[arrow] (load) -- (shape);
\draw[arrow] (shape) -- (missing);
\draw[arrow] (missing) -- (desc);
\draw[arrow] (desc.east) -- ++(1.25,0) |- (uni.west);
\draw[arrow] (uni) -- (bi);
\draw[arrow] (bi) -- (multi);
\draw[arrow] (multi) -- (concl);
\end{tikzpicture}
\end{center}

% ============================================================
\section{Case Study: The Titanic Dataset}
\index{Titanic}

\subsection{Loading and First Inspection}

\begin{codeblock}[title=\textbf{Python}]
\begin{minted}{python}
import pandas as pd
import numpy as np
import seaborn as sns
import matplotlib.pyplot as plt

titanic = sns.load_dataset('titanic')
print(f"Dimensions : {titanic.shape}")
print(titanic.dtypes)
\end{minted}
\end{codeblock}

\begin{codeoutput}
\begin{minted}{text}
Dimensions : (891, 15)
survived         int64
pclass           int64
sex             object
age            float64
sibsp            int64
parch            int64
fare           float64
embarked        object
class         category
who             object
adult_male        bool
deck          category
embark_town     object
alive           object
alone             bool
\end{minted}
\end{codeoutput}

\subsection{Missing Values}

\begin{codeblock}[title=\textbf{Python}]
\begin{minted}{python}
missing = titanic.isnull().sum()
missing_pct = (missing / len(titanic) * 100).round(1)
missing_df = pd.DataFrame({'Missing': missing,
                           'Percentage': missing_pct})
print(missing_df[missing_df['Missing'] > 0])
\end{minted}
\end{codeblock}

\begin{codeoutput}
\begin{minted}{text}
             Missing  Percentage
age              177        19.9
embarked           2         0.2
deck             688        77.2
embark_town        2         0.2
\end{minted}
\end{codeoutput}

\begin{remark}
The variable \py{deck} has 77\,\% missing values: it will likely be unusable without heavy imputation. The variable \py{age} has about 20\,\% missing, which is manageable.
\end{remark}

\subsection{Descriptive Statistics}

\begin{codeblock}[title=\textbf{Python}]
\begin{minted}{python}
print(titanic.describe())
print("\n--- Categorical Variables ---")
print(titanic[['sex', 'embarked', 'class']].describe())
\end{minted}
\end{codeblock}

\begin{codeoutput}
\begin{minted}{text}
         survived     pclass        age      sibsp      parch       fare
count  891.000000  891.000000  714.00000  891.00000  891.00000  891.00000
mean     0.383838    2.308642   29.69912    0.52301    0.38159   32.20421
std      0.486592    0.836071   14.52650    1.10274    0.80606   49.69343
min      0.000000    1.000000    0.42000    0.00000    0.00000    0.00000
25%      0.000000    2.000000   20.12500    0.00000    0.00000    7.91040
50%      0.000000    3.000000   28.00000    0.00000    0.00000   14.45420
75%      1.000000    3.000000   38.00000    1.00000    0.00000   31.00000
max      1.000000    3.000000   80.00000    8.00000    6.00000  512.32920

--- Categorical Variables ---
          sex embarked class
count     891      889   891
unique      2        3     3
top      male        S Third
freq      577      644   491
\end{minted}
\end{codeoutput}

% ============================================================
\section{Univariate Analysis}
\index{univariate analysis}

\begin{codeblock}[title=\textbf{Python}]
\begin{minted}{python}
fig, axes = plt.subplots(2, 2, figsize=(12, 9))

sns.histplot(titanic['age'].dropna(), bins=30, kde=True,
             ax=axes[0,0], color='steelblue')
axes[0,0].set_title('Age Distribution')

sns.countplot(data=titanic, x='pclass', ax=axes[0,1],
              palette='viridis')
axes[0,1].set_title('Distribution by Class')

sns.countplot(data=titanic, x='survived', ax=axes[1,0],
              palette='Set2')
axes[1,0].set_xticklabels(['Deceased', 'Survived'])
axes[1,0].set_title('Survival')

sns.histplot(titanic['fare'], bins=40, ax=axes[1,1],
             color='coral')
axes[1,1].set_title('Fare Distribution')

plt.tight_layout()
plt.show()
\end{minted}
\end{codeblock}

\begin{codeoutput}
Four plots: age approximately follows a normal distribution centered around 29 years; third class is the most represented (491); 62\,\% of passengers did not survive; the fare is highly skewed with a long right tail.
\end{codeoutput}

% ============================================================
\section{Bivariate Analysis}
\index{bivariate analysis}

\subsection{Survival by Sex and by Class}

\begin{codeblock}[title=\textbf{Python}]
\begin{minted}{python}
fig, axes = plt.subplots(1, 3, figsize=(15, 5))

# Survival by sex
sns.barplot(data=titanic, x='sex', y='survived',
            ax=axes[0], palette='coolwarm', ci=95)
axes[0].set_title('Survival Rate by Sex')
axes[0].set_ylabel('Survival Rate')

# Survival by class
sns.barplot(data=titanic, x='pclass', y='survived',
            ax=axes[1], palette='viridis', ci=95)
axes[1].set_title('Survival Rate by Class')

# Survival by sex and class
sns.barplot(data=titanic, x='pclass', y='survived',
            hue='sex', ax=axes[2], palette='Set1', ci=95)
axes[2].set_title('Survival by Class and Sex')

plt.tight_layout()
plt.show()
\end{minted}
\end{codeblock}

\begin{codeoutput}
Women have a survival rate of 74\,\% compared to 19\,\% for men. First class has 63\,\% survival, second class 47\,\%, and third class 24\,\%. The combination shows that first-class women survived at nearly 97\,\%.
\end{codeoutput}

\subsection{Correlations and Contingency Table}
\index{contingency table}

\begin{codeblock}[title=\textbf{Python}]
\begin{minted}{python}
# Contingency table
ct = pd.crosstab(titanic['sex'], titanic['survived'],
                 margins=True)
ct.columns = ['Deceased', 'Survived', 'Total']
print(ct)

# Numerical correlation
cols_num = ['survived', 'pclass', 'age', 'sibsp', 'parch', 'fare']
corr = titanic[cols_num].corr()
print("\nCorrelation Matrix:")
print(corr.round(2))
\end{minted}
\end{codeblock}

\begin{codeoutput}
\begin{minted}{text}
         Deceased  Survived  Total
sex
female       81        233    314
male        468        109    577
Total       549        342    891

Correlation Matrix:
          survived  pclass    age  sibsp  parch   fare
survived      1.00   -0.34  -0.08  -0.04   0.08   0.26
pclass       -0.34    1.00  -0.37   0.08   0.02  -0.55
age          -0.08   -0.37   1.00  -0.31  -0.19   0.10
fare          0.26   -0.55   0.10  -0.24   0.22   1.00
\end{minted}
\end{codeoutput}

% ============================================================
\section{Multivariate Analysis}
\index{multivariate analysis}

\begin{codeblock}[title=\textbf{Python}]
\begin{minted}{python}
fig, ax = plt.subplots(figsize=(8, 6))
sns.heatmap(corr, annot=True, fmt='.2f', cmap='RdBu_r',
            center=0, ax=ax, square=True,
            linewidths=0.5)
ax.set_title('Correlation Heatmap -- Titanic')
plt.tight_layout()
plt.show()
\end{minted}
\end{codeblock}

\begin{codeoutput}
Heatmap showing correlations: survival is negatively correlated with class ($-0.34$) and positively with fare ($0.26$). Class and fare are strongly anti-correlated ($-0.55$).
\end{codeoutput}

% ============================================================
\section{Outlier Detection}
\index{outliers}\index{IQR}

\begin{definition}[IQR Method]
An observation is considered an outlier if it falls outside the interval:
\[
\bigl[Q_1 - 1.5 \times \text{IQR},\; Q_3 + 1.5 \times \text{IQR}\bigr]
\]
where $\text{IQR} = Q_3 - Q_1$ is the interquartile range\index{interquartile range}.
\end{definition}

\begin{codeblock}[title=\textbf{Python}]
\begin{minted}{python}
def detect_outliers_iqr(df, column):
    Q1 = df[column].quantile(0.25)
    Q3 = df[column].quantile(0.75)
    IQR = Q3 - Q1
    lower_bound = Q1 - 1.5 * IQR
    upper_bound = Q3 + 1.5 * IQR
    outliers = df[(df[column] < lower_bound) |
                  (df[column] > upper_bound)]
    return outliers, lower_bound, upper_bound

outliers_fare, b_low, b_high = detect_outliers_iqr(titanic, 'fare')
print(f"Bounds: [{b_low:.2f}, {b_high:.2f}]")
print(f"Number of outliers (fare): {len(outliers_fare)}")
print(f"Max fare: {titanic['fare'].max():.2f}")
\end{minted}
\end{codeblock}

\begin{codeoutput}
\begin{minted}{text}
Bounds: [-26.72, 65.63]
Number of outliers (fare): 116
Max fare: 512.33
\end{minted}
\end{codeoutput}

\begin{remark}
116 passengers (13\,\%) have a fare considered as an outlier by the IQR method. These are mainly first-class passengers. They should not be automatically removed: these are legitimate values that reflect the data structure.
\end{remark}

% ============================================================
\section{Exercises}

\begin{exercise}[$\star$]
Load the \py{titanic} dataset and display the first 5 rows, the dimensions, the column types, and the number of missing values per column.
\end{exercise}

\begin{exercise}[$\star$]
Compute the overall survival rate, then the survival rate by port of embarkation (\py{embarked}). Which port letter has the highest survival rate?
\end{exercise}

\begin{exercise}[$\star\star$]
Create a chart with \py{sns.catplot} of type \py{'bar'} showing the survival rate by class and by sex. Add an informative title. Interpret the results in terms of evacuation priority.
\end{exercise}

\begin{exercise}[$\star\star$]
Create a histogram of passenger age, separated by survival (\py{hue='survived'}). Are survivors younger on average? Verify with \py{groupby('survived')['age'].mean()}.
\end{exercise}

\begin{exercise}[$\star\star\star$]
Apply the IQR method to detect outliers in the \py{age} variable. How many are there? Create an annotated box plot showing the IQR bounds and the outliers. Discuss the relevance of removing these observations.
\end{exercise}

\begin{exercise}[$\star\star\star$]
Perform a complete EDA of the \py{penguins} dataset (\py{sns.load\_dataset('penguins')}). Your analysis should include: (a) general inspection, (b) missing value handling, (c) univariate distributions, (d) bivariate analysis by species, (e) correlation heatmap. Conclude by formulating 3 testable hypotheses.
\end{exercise}

% ============================================================
\begin{keyformulas}
\textbf{Chapter Summary -- Exploratory Data Analysis}
\begin{itemize}
    \item \textbf{Inspection}: \py{df.shape}, \py{df.dtypes}, \py{df.info()}, \py{df.head()}.
    \item \textbf{Missing values}: \py{df.isnull().sum()}, \py{df.isnull().mean() * 100}.
    \item \textbf{Descriptive}: \py{df.describe()}, \py{df['col'].value\_counts()}.
    \item \textbf{Univariate}: \py{sns.histplot()}, \py{sns.countplot()}, \py{sns.boxplot()}.
    \item \textbf{Bivariate}: \py{sns.barplot(x, y)}, \py{sns.scatterplot()}, \py{pd.crosstab()}.
    \item \textbf{Multivariate}: \py{df.corr()}, \py{sns.heatmap()}, \py{sns.pairplot()}.
    \item \textbf{Outliers (IQR)}: outlier if $x < Q_1 - 1.5 \cdot \text{IQR}$ or $x > Q_3 + 1.5 \cdot \text{IQR}$.
    \item \textbf{Tukey's philosophy}: explore first, model later.
\end{itemize}
\end{keyformulas}
